\documentclass[12pt]{article}

% DEFAULT PACKAGE SETUP
\usepackage{setspace,graphicx,epstopdf,amsmath,amsfonts,amssymb,amsthm,versionPO,threeparttable}
\usepackage{marginnote,datetime,enumitem,subfig,rotating,fancyvrb,tabularx,booktabs}
\usepackage{float}
% \usepackage[longnamesfirst]{natbib}
\usepackage[]{natbib}
\newcolumntype{Y}{>{\raggedleft\arraybackslash}X}% raggedleft column X
\interfootnotelinepenalty=10000 %% Completely prevent breaking of footnotes

% Define \numberthis
\newcommand{\numberthis}{\addtocounter{equation}{1}\tag{\theequation}}

\usdate
\usepackage{lscape}
% \usepackage{endfloat}
% \usepackage[nolists]{endfloat}
\usepackage{rotating} 

\usepackage[justification=justified]{caption}
% \usepackage{subcaption}
% These next lines allow including or excluding different versions of text
% using versionPO.sty

\excludeversion{notes}  % Include notes?
\includeversion{links}  % Turn hyperlinks on?

\usepackage{eurosym}
% Turn off hyperlinking if links is excluded
\iflinks{}{\hypersetup{draft=true}}

% Notes options
\ifnotes{%
 \usepackage[margin=1in,paperwidth=10in,right=2.5in]{geometry}%
 \usepackage[textwidth=1.4in,shadow,colorinlistoftodos]{todonotes}%
}{%
 \usepackage[margin=1in]{geometry}%
 \usepackage[disable]{todonotes}%
}

\DeclareMathOperator*{\argmin}{arg\,min}


% Allow todonotes inside footnotes without blowing up LaTeX
% Next command works but now notes can overlap. Instead, we'll define
% a special footnote note command that performs this redefinition.
% \renewcommand{\marginpar}{\marginnote}%

% Save original definition of \marginpar
\let\oldmarginpar\marginpar

% Workaround for todonotes problem with natbib (To Do list title comes out wrong)
\makeatletter\let\chapter\@undefined\makeatother % Undefine \chapter for todonotes

% Define note commands
\newcommand{\smalltodo}[2][] {\todo[caption={#2}, size=\scriptsize, fancyline, #1] {\begin{spacing}{.5}#2\end{spacing}}}
\newcommand{\rhs}[2][]{\smalltodo[color=green!30,#1]{{\bf RS:} #2}}
\newcommand{\rhsnolist}[2][]{\smalltodo[nolist,color=green!30,#1]{{\bf RS:} #2}}
\newcommand{\rhsfn}[2][]{% To be used in footnotes (and in floats)
 \renewcommand{\marginpar}{\marginnote}%
 \smalltodo[color=green!30,#1]{{\bf RS:} #2}%
 \renewcommand{\marginpar}{\oldmarginpar}}
% \newcommand{\textnote}[1]{\ifnotes{{\noindent\color{red}#1}}{}}
\newcommand{\textnote}[1]{\ifnotes{{\colorbox{yellow}{{\color{red}#1}}}}{}}

% Command to start a new page, starting on odd-numbered page if twoside option
% is selected above
\newcommand{\clearRHS}{\clearpage\thispagestyle{empty}\cleardoublepage\thispagestyle{plain}}

% Number paragraphs and subparagraphs and include them in TOC
\setcounter{tocdepth}{2}

% JF-specific includes:

\usepackage{indentfirst} % Indent first sentence of a new section.
\usepackage{endnotes} % Use endnotes instead of footnotes
\usepackage{rfs}  % JF-specific formatting of sections, etc.
% \usepackage[labelfont=bf,labelsep=period]{caption} % Format figure captions
\captionsetup[table]{labelsep=none}

\usepackage{color}%Color
% Define theorem-like commands and a few random function names.
\newtheorem{condition}{CONDITION}
\newtheorem{corollary}{COROLLARY}
\newtheorem{proposition}{PROPOSITION}
\newtheorem{obs}{OBSERVATION}
\newtheorem{lem}{LEMMA}
\newtheorem{theorem}{THEOREM}
\newtheorem{assumption}{ASSUMPTION}
\newcommand{\argmax}{\mathop{\rm arg\,max}}
\newcommand{\sign}{\mathop{\rm sign}}
\newcommand{\defeq}{\stackrel{\rm def}{=}}

\usepackage{multirow}
\usepackage{booktabs}


\usepackage{amsmath}
\usepackage{amsthm}
\usepackage{amsfonts}
\usepackage{relsize}
\usepackage{tikz}
\usepackage{pgfplotstable} 
\usepgfplotslibrary{dateplot}
\usetikzlibrary{fit}
\usepackage{graphicx}
\usepackage{eurosym}
\usepackage{standalone}
% \usepackage{filecontents}

\usepackage[section]{placeins} % keep figures in the same section


\usepackage{appendix}
\appendixtitleon
\newcommand\independent{\protect\mathpalette{\protect\independenT}{\perp}}
\def\independenT#1#2{\mathrel{\rlap{$#1#2$}\mkern2mu{#1#2}}}

%\onehalfspacing
%\linespread{1.25}

\setstretch{1.5}
%\doublespacing

\definecolor{zahi}{RGB}{255,0,0}
\newcommand{\zahi}[1]{\textcolor{zahi}{[zahi: #1]}}



%%%%% Pgfplots %%%%%
\usepackage{pgfplots}
\usepgfplotslibrary{groupplots}

\pgfplotsset{compat=newest}
\usetikzlibrary{patterns}

%%%% Table caption package %%%%%
\usepackage{ragged2e}
\usepackage{longtable}

%%%% Figure imports %%%%
\usepackage{import}

%%%% Color cells in tables %%%%
\usepackage{tabulary,colortbl}
\newcolumntype{C}[1]{>{\centering\let\newline\\\arraybackslash\hspace{0pt}}m{#1}}


\definecolor{andrea}{RGB}{255,0,0}
\newcommand{\andrea}[1]{\textcolor{andrea}{[andrea: #1]}}

\definecolor{yang}{RGB}{255,0,0}
\newcommand{\yang}[1]{\textcolor{yang}{[yang: #1]}}
\newcommand{\jay}[1]{\textcolor{yang}{[j: #1]}}


\usepackage[colorlinks=true,linkcolor=red,anchorcolor=black,citecolor=black, filecolor=black,menucolor=black,runcolor=black,urlcolor=black]{hyperref}
\usepackage{parskip}



%%%%%%%%%%%%%%%%% For plots: get min, max etc


\newcommand*{\GetElement}[4]{%
    \pgfplotstablegetelem{#2}{#3}\of{#1}%
    \let#4\pgfplotsretval
}

% get the first element of a column and store it in a macro
% #1: table
% #2: column (name or [index]<index>)
% #3: macro for value
\newcommand*{\GetFirstElement}[3]{%
    \GetElement{#1}{0}{#2}{#3}%
}

% get the last element of a column and store it in a macro
% #1: table
% #2: column (name or [index]<index>)
% #3: macro for value
\newcommand*{\GetLastElement}[3]{%
    \pgfplotstablegetrowsof{#1}%
    \pgfmathsetmacro\lastrowidx{int(\pgfplotsretval - 1)}%
    \GetElement{#1}{\lastrowidx}{#2}{#3}%
}

% get maximum and minimum value of a column and store them in macros
% #1: table
% #2: column
% #3: macro for max
% #4: macro for min
\newcommand*{\GetMinMax}[4]{%
    \pgfmathsetmacro#3{-16383}%
    \pgfmathsetmacro#4{16383}%
    \pgfplotstableforeachcolumnelement{#2}\of{#1}\as\cellValue{%
        \ifx\cellValue\@empty\else
            \pgfmathsetmacro{#3}{max(#3,\cellValue)}%
            \pgfmathsetmacro{#4}{min(#4,\cellValue)}%
        \fi
    }
}

% get average of column and store it in a macro
% #1: table
% #2: column
% #3: macro for average
\newcommand*{\GetAverage}[3]{%
    \pgfplotstablegetrowsof{#1}%
    \let\rownumber\pgfplotsretval
    \def#3{0}%
    \pgfplotstableforeachcolumnelement{#2}\of{#1}\as\cellValue{%
        \ifx\cellValue\@empty
            % don't count emplty cells
            \pgfmathsetmacro{\rownumber}{\rownumber - 1}%
        \else
            \pgfmathsetmacro{#3}{#3 + \cellValue}%
        \fi
    }
    \pgfmathsetmacro{#3}{#3 / \rownumber}%
}

%%%%%%%%%%%%%%%%%%%%%%%%%%%%%%%%%%%%%%%%%%%%%%%%%%%%%%

%%%% Quintile colors (Green to Red) %%%% 
\definecolor{q1}{rgb}{ .973,  .412,  .42}
\definecolor{q2}{rgb}{ .992,  .749,  .486}
\definecolor{q3}{rgb}{ 1,  .922,  .518}
\definecolor{q4}{rgb}{ .678,  .827,  .498}
\definecolor{q5}{rgb}{ .388,  .745,  .482}
%%%%%%%%%%%%%%%%%%%%%%%%%%%%%%%%%%%%%%%%

% optional: end floats
\usepackage[nolists]{endfloat}

\begin{document}
\setlist{noitemsep}  % Reduce space between list items (itemize, enumerate, etc.)
\onehalfspacing      % Use 1.5 spacing
\renewcommand{\thempfootnote}{\fnsymbol{mpfootnote}}


\vspace{1.5cm}
 \author{
% {\large Jiacui Li\hspace{0.00in}}
 }

\title{\vspace*{0.0in}\Large \bf 
Considerations in government intervention in stock markets
% \thanks{\rm  Li (\url{jiacui.li@eccles.utah.edu}) is at David Eccles School of Business, University of Utah. }
}

% \vspace{0.2cm}

\date{\today}              % No date for final submission

% Create title page with no page number

\maketitle


\thispagestyle{empty}
\onehalfspacing
\centerline{\bf ABSTRACT}

\noindent I extended the \citet{brunnermeier2022china} (BSX) framework to incorporate two government policy tools: (1) direct trading against noise traders, and (2) adjusting financing conditions to influence investor demand. It also models transparency through a noisy public signal about future government interventions.



Here are the key conclusions:
\begin{itemize}
    \item Greater transparency unambiguously reduces return volatility, but it lowers government trading profits because investors front-run anticipated interventions.

    \item In the model, direct intervention and financing policy are equally effective in stabilizing prices. Their real-world differences stem from government costs, with financing policy being ``cheaper'' (does not require fiscal outlays). 
\end{itemize}

\medskip

\vspace{0.1cm}

% \noindent \textit{JEL classification:} G11, G12, G14\\
% \noindent \textit{Keywords:} Lead-Lag, Cross-Stock Momentum, Factor Momentum, Reversal, Underreaction

\thispagestyle{empty}


\setcounter{page}{1}


%%%%%%%%%%%%%%%%%%%%%
\newpage
\section{Model}

\subsection{Set up}

Time is discrete, \(t = 0,1,2,\dots\).

\paragraph{Assets.}
There is a risky asset with one share outstanding. Dividends follow
\begin{align*}
    D_t & = V_t + \sigma_D \epsilon_t^D, \qquad \epsilon_t^D \stackrel{\text{i.i.d.}}{\sim} N(0,1),
\end{align*}
where the fundamental component \(V_t\) is persistent:
\begin{align*}
    V_t & = \rho_V V_{t-1} + \sigma_V \epsilon_t^V, \qquad \epsilon_t^V \stackrel{\text{i.i.d.}}{\sim} N(0,1).
\end{align*}

The risk‑free rate \(R^f > 1\) is set exogenously. The excess return on one share of the risky asset is
\begin{align}
    R_{t+1} & = D_{t+1} + P_{t+1} - R^f P_t . \label{excess_return}
\end{align}


\paragraph{Financing cost.}
Investors (introduced below) incur a per‑share holding cost \(\phi_t\).  
Thus the net excess return from holding one share is \(R_{t+1} - \phi_t\).  
The cost \(\phi_t\) can be interpreted as a reduced‑form way of capturing leverage or margin‑financing costs.\footnote{A more serious model of margin constraints will likely break the linear‑Gaussian structure of this class of models.}

\paragraph{Agents.}

\begin{itemize}
\item \underline{Noise traders.}
Noise traders submit net purchases
\[
N_t = \sigma_N \epsilon_t^N, \qquad \epsilon_t^N \stackrel{\text{i.i.d.}}{\sim} N(0,1).
\]

\item \underline{Investors.}
A unit‑mass continuum of myopic investors have constant absolute risk aversion (CARA) utility with absolute risk aversion \(\gamma\). Before submitting orders at time $t$, their information set is  
\[
\mathcal{F}_t = \{\, V_{s+1}, N_s, D_s, \phi_t \,\}_{s \le t}.
\]  
Because they incur the per‑share holding cost \(\phi_t\), their demand for the risky asset is
\begin{align}
    X_t^I = \frac{E_t(R_{t+1}) - \phi_t}{\gamma \,\mathrm{Var}_t(R_{t+1})}. \label{demand_with_phi}
\end{align}
\end{itemize}

\paragraph{Government.}
Following BSX, the government aims to reduce return volatility by leaning against noise traders. It observes \(N_t\) with noise:\footnote{In BSX's specification, the noise in the government's signal has the same standard deviation as \(N_t\) (\(\sigma_N\)). This normalization reduces the number of free parameters; I think the qualitative insights would carry over if the noise were parameterized separately, e.g., as \(\sigma_G\).}
\begin{align}
    \hat{N}_t^G = N_t - \sigma_N G_t, \qquad G_t \stackrel{\text{i.i.d.}}{\sim} N(0,1). \label{govt_noisy_signal}
\end{align}


The government has two policy instruments:

\begin{itemize}
    \item \underline{Direct trading.} The government submits the order
    \begin{align}
        X_t^G = - \underbrace{\vartheta_N}_{\text{intensity}} \cdot \underbrace{\bigl(N_t - \sigma_N G_t\bigr)}_{\hat{N}_t^G}. \label{govt_demand}
    \end{align}

    \item \underline{Adjusting financing cost.} The government can set the per‑share holding cost \(\phi_t\) faced by investors. Lowering (raising) \(\phi_t\) encourages (discourages) investor demand. The policy rule is~\footnote{Setting \(\mathbb{E}[\phi_t]=0\) is essentially without loss of generality: a constant positive mean for \(\phi_t\) would act like an increase in the effective risk‑free rate \(R^f\).}
    \begin{align}
        \phi_t = \underbrace{\theta_N}_{\text{intensity}} \cdot \underbrace{\bigl(N_t - \sigma_N G_t\bigr)}_{\hat{N}_t^G}. \label{phi}
    \end{align}
    
\end{itemize}

\underline{Transparency.}
The government can also choose how transparent it is about its future actions. Before investors submit orders at time \(t\), the government releases a public signal about next period's shock \(G_{t+1}\):
\begin{align*}
    g_t = G_{t+1} + \tau_G^{-1/2} \epsilon_t^G, \qquad \epsilon_t^G \stackrel{\text{i.i.d.}}{\sim} N(0,1).
\end{align*}
Higher precision \(\tau_G\) corresponds to greater transparency. 

\subsection{Solving for equilibrium given govenment policy}

We first take the government policy parameters \((\vartheta_N, \theta_N, \tau_G)\) as given and solve for the rational‑expectations equilibrium.

\paragraph{Learning from the government's signal.}
Investors form their expectation of \(G_{t+1}\) using the public signal \(g_t\):
\begin{align*}
    \hat{G}_{t+1} &\equiv \mathbb{E}[\,G_{t+1} \mid g_t\,] 
                  = \frac{\operatorname{Cov}(G_{t+1}, g_t)}{\operatorname{Var}(g_t)} \, g_t 
                  = \frac{1}{1 + \tau_G^{-1}} \, g_t .
\end{align*}
The residual uncertainties about government shocks one and two periods ahead are
\begin{align*}
    \operatorname{Var}_t(G_{t+1}) 
        &= \operatorname{Var}(G_{t+1}) - \operatorname{Var}(\hat{G}_{t+1}) 
          = 1 - \frac{1}{1+\tau_G^{-1}} 
          = \frac{\tau_G^{-1}}{1+\tau_G^{-1}},\\[4pt]
    \operatorname{Var}_t(\hat{G}_{t+2}) 
        &= \operatorname{Var}_t\!\left(\frac{1}{1+\tau_G^{-1}} \, g_{t+1}\right) \\
        &= \operatorname{Var}_t\!\left(
              \frac{1}{1+\tau_G^{-1}} \,
              \bigl(G_{t+2} + \tau_G^{-1/2} \epsilon_{t+1}^G\bigr)
           \right) 
          = \frac{1}{1+\tau_G^{-1}} .
\end{align*}

\paragraph{Price conjecture.}
We posit a linear equilibrium price of the form
\begin{align}
    P_t = \underbrace{\bar{p}}_{\text{risk discount}} 
          + \underbrace{\frac{1}{R^f - \rho_V} V_{t+1}}_{\text{fundamental}} 
          + \underbrace{p_G G_t}_{\text{current govt shock}} 
          + \underbrace{p_{\hat{G}} \hat{G}_{t+1}}_{\text{expected govt shock}} 
          + \underbrace{p_N N_t}_{\text{noise trading}}. \label{price_conjecture}
\end{align}
Intuitively, stronger government intervention reduces the coefficient on noise trading (\(p_N\)) but increases the coefficients on current and expected government shocks (\(p_G\) and \(p_{\hat{G}}\)). 

\paragraph{Conditional moments of returns.}
Substituting the conjecture into \eqref{excess_return} gives
\begin{align*}
    \mathbb{E}_t[R_{t+1}] 
        &= \mathbb{E}_t[D_{t+1}] + \mathbb{E}_t[P_{t+1}] - R^f P_t \\
        &= \frac{R^f}{R^f - \rho_V} V_{t+1} + p_G \hat{G}_{t+1} + \bar{p} - R^f P_t, \\[6pt]
    \operatorname{Var}_t(R_{t+1}) 
        &= \operatorname{Var}_t\!\bigl(D_{t+1} + P_{t+1}\bigr) \\
        &= \operatorname{Var}_t\!\Bigl(
                \sigma_D \epsilon_{t+1}^D 
                + \frac{1}{R^f - \rho_V} V_{t+2} 
                + p_G G_{t+1} 
                + p_{\hat{G}} \hat{G}_{t+2} 
                + p_N N_{t+1}
           \Bigr) \\
        &= \underbrace{\sigma_D^2 
                + \left(\frac{1}{R^f - \rho_V}\right)^2 \sigma_V^2}_{\text{fundamental risk}}
           + \underbrace{p_G^2 \frac{\tau_G^{-1}}{1+\tau_G^{-1}}}_{\text{policy risk}}
           + \underbrace{p_{\hat{G}}^2 \frac{1}{1+\tau_G^{-1}}}_{\text{policy signal risk}}
           + \underbrace{p_N^2 \sigma_N^2}_{\text{noise trader risk}} .
           \numberthis \label{var_R}
\end{align*}

\paragraph{Investor demand and market clearing.}
Using \eqref{demand_with_phi}, investor demand is
\[
X_t^I = \frac{\mathbb{E}_t[R_{t+1}] - \phi_t}
            {\gamma \operatorname{Var}_t(R_{t+1})}
      = \frac{1}{\gamma \operatorname{Var}_t(R_{t+1})}
        \left(
            \frac{R^f}{R^f - \rho_V} V_{t+1} 
            + p_G \hat{G}_{t+1} 
            + \bar{p} 
            - R^f P_t 
            - \phi_t
        \right).
\]
Market clearing requires
\begin{align*}
    X_t^I + X_t^G + N_t = 1 .
\end{align*}
Substituting the expressions for \(X_t^I\) and \(X_t^G\) yields
\begin{align*}
    \frac{1}{\gamma \operatorname{Var}_t(R_{t+1})}
    \left(
        \frac{R^f}{R^f - \rho_V} V_{t+1} 
        + p_G \hat{G}_{t+1} 
        + \bar{p} 
        - R^f P_t 
        - \phi_t
    \right)
    + (1-\vartheta_N) N_t 
    + \vartheta_N \sigma_N G_t = 1 .
\end{align*}

\paragraph{Solving for the price.}
Solving the market‑clearing equation for \(P_t\) gives
\begin{align}
    P_t = \frac{1}{R^f}
          \Bigl\{
              \gamma \operatorname{Var}_t(R_{t+1})
                  \bigl[ (1-\vartheta_N) N_t + \vartheta_N \sigma_N G_t - 1 \bigr]
              + \frac{R^f}{R^f - \rho_V} V_{t+1}
              + p_G \hat{G}_{t+1}
              + \bar{p}
              - \phi_t
          \Bigr\}. \label{price_market_clearing_with_phi}
\end{align}

\paragraph{Matching coefficients.}
Comparing \eqref{price_market_clearing_with_phi} with the conjectured form \eqref{price_conjecture} and using \(\phi_t = \theta_N (N_t - \sigma_N G_t)\) from \eqref{phi} yields the equilibrium coefficients:
\begin{align}
    \bar{p} &= -\,\frac{\gamma \operatorname{Var}_t(R_{t+1})}{R^f - 1}, \label{pbar}\\[4pt]
    p_G     &= \frac{\bigl[\gamma \operatorname{Var}_t(R_{t+1}) \vartheta_N + \theta_N\bigr] \sigma_N}{R^f}, \label{pG}\\[4pt]
    p_{\hat{G}} &= \frac{p_G}{R^f}, \label{pGhat}\\[4pt]
    p_N     &= \frac{\gamma \operatorname{Var}_t(R_{t+1}) (1-\vartheta_N) - \theta_N}{R^f}. \label{pN}
\end{align}
These expressions verify that the linear conjecture \eqref{price_conjecture} indeed constitutes an equilibrium.


\paragraph{Discussion of the equilibrium.}

Holding \(\operatorname{Var}_t(R_{t+1})\) fixed, the coefficients in \eqref{pbar}--\eqref{pN} have intuitive interpretations.  
The constant \(\bar{p}\) is a risk‑based discount.  
Increasing either direct intervention (\(\vartheta_N\)) or margin‑based intervention (\(\theta_N\)) raises the coefficients on current and expected government shocks (\(p_G\) and \(p_{\hat{G}}\)) while reducing the coefficient on noise trading (\(p_N\)).  

Of course, \(\operatorname{Var}_t(R_{t+1})\) is endogenous.  Inspecting \eqref{var_R} shows that greater transparency (higher \(\tau_G\)) lowers the conditional variance of returns.\footnote{This follows because \(p_{\hat{G}} = p_G/R^f < p_G\); the reduction in policy risk outweighs the increase in policy‑signal risk.}

\paragraph{Government trading profits.}

Because the government leans against noise traders, as long as it does not over-trade, its trades are profitable on average.  Its expected per‑period trading profit (in excess of the risk‑free rate) is
\begin{align*}
    \mathbb{E}[\Pi_{t+1}] &\equiv \mathbb{E}\!\bigl[ X_t^G \cdot R_{t+1} \bigr] 
                          = \mathbb{E}\!\bigl[ X_t^G \cdot \mathbb{E}_t[R_{t+1}] \bigr] \\
    &= \mathbb{E}\Bigl\{ X_t^G \cdot 
        \bigl[ \phi_t + \gamma \operatorname{Var}_t(R_{t+1}) (1 - N_t - X_t^G) \bigr] \Bigr\} \\
    &= \mathbb{E}\bigl[ X_t^G \phi_t \bigr] 
       + \gamma \operatorname{Var}_t(R_{t+1}) \,
         \mathbb{E}\!\bigl[ X_t^G (1 - N_t - X_t^G) \bigr].
\end{align*}
Using the policy rules \eqref{govt_demand} and \eqref{phi}, we obtain
\begin{align}
    \mathbb{E}[\Pi_{t+1}] 
    &= - \underbrace{2\vartheta_N\theta_N\sigma_N^2}_{\text{policy conflict}}
       + \underbrace{\gamma \operatorname{Var}_t(R_{t+1}) \,
                     \vartheta_N (1-2\vartheta_N)\,\sigma_N^2}_{\text{trading gains}} .
    \label{govt_profit}
\end{align}

The first term reflects a \emph{policy conflict} that arises only when both tools are active.  
Adjusting the financing cost (\(\theta_N > 0\)) diminishes the profitability of the government's own trades.%
\footnote{For example, when \(N_t<0\) the government buys shares (\(X_t^G>0\)) but also lowers \(\phi_t\), which reduces the expected excess return investors require and thereby lowers the price appreciation the government can capture.}

The second term captures the pure trading gains from leaning against noise traders.  
It is positive as long as \(\vartheta_N < 1/2\)—the natural case when the government aims to reduce volatility rather than amplify it.  

Crucially, government profits depend on transparency \(\tau_G\) through \(\operatorname{Var}_t(R_{t+1})\).  
Increasing transparency (\(\tau_G \uparrow\)) lowers \(\operatorname{Var}_t(R_{t+1})\) and therefore reduces trading profits.  
This is the \emph{front‑running} channel: a more predictable government is partially anticipated by investors, eroding its ability to profit from supplying liquidity against noise.


\subsection{Implications for optimal policy}

The analysis reveals a trade‑off between volatility reduction and government budgetary considerations—the latter are not explicitly modeled here.

\begin{itemize}
    \item \textbf{Transparency.}  
    For given intervention intensities \((\vartheta_N,\theta_N)\), higher transparency (\(\tau_G\)) always lowers return volatility.  
    Within the model, the sole cost of transparency is a reduction in government trading profits.  

    \begin{itemize}
        \item The front‑running effect arises because the government commits to a \emph{quantity‑based} intervention rule. If instead the government could credibly announce a \emph{price‑based} target (e.g., ``I will keep the price above a certain level''), market participants might move the price to that level without the government needing to trade at all.
    \end{itemize}

    \item \textbf{Direct trading versus margin policy.}  
    From a pure volatility‑reduction standpoint, the two tools are quite similar.  
    For any given direct‑trading intensity \(\vartheta_N\), there exists a margin‑policy intensity \(\theta_N\) that produces the same equilibrium coefficients \(p_G, p_N\) and hence the same price impact on noise trading.\footnote{This exact substitutability is a consequence of the linear‑Gaussian structure and may not generalize to richer settings.}

    \begin{itemize}
        \item The key practical distinction likely lies in \emph{implementation costs}.  
    Direct trading exposes the government’s balance sheet to market risk and requires substantial liquidity.  
    Adjusting margin requirements, in contrast, involves minimal fiscal outlay and no direct market position.  
    If the government faces risk or position limits, it would naturally prefer to induce private investors to absorb shocks rather than doing so itself.
    \end{itemize}
\end{itemize}

\subsection{J - thoughts about step-by-step extensions}

1) Perhaps just add a ``risk aversion'', interpreted as ``ammo'', to the NT. And see what happens. 

2) Is there any sense in which we should model the NT as optimizing over longer horizons? 



\newpage
\section{New thoughts -- assuming objective is reducing volatility}

There is a huge literature on central banks that touch on optimal intervention, degree of opacity, and credibility issues that we can build on. They often have some inflation target, and central bank tweaks interest rate or money supply to hit that target. 

How are things different here? A few things:
\begin{itemize}
    \item CB's objective is to reduce deviation of inflation from an \textit{observable} target (e.g. 2\%). 
    \begin{itemize}
        \item In contrast, NT (national team) wants to reduce price deviation from ``fundamentals'' --- a target that is not seen, and arguably some investors know more about. Here, we have to take seriously the idea that NT has a rather noisy signal of the ``right target''. 
        \item To the extent that prices eventually tend towards fundamentals, NT cannot distort prices forever. In a sense, all intervention are ``temporary'' in nature. 
    \end{itemize}

    \item NT and CB both have ``credibility'' issues, but for different reasons. In CB, people are concerned about time-inconsistency (Kydland \& Prescott (1977) and Barro \& Gordon (1983)). Or, less importantly, as in Backus \& Driffill (1985), CB has incentive to ``pretend to be more serious than it is'' to make its intervention more potent. 
    \begin{itemize}
        \item For NT, it is more about ``does the NT have enough amo'', while CB is often assumed to have infinite amo. 
        \item Of course, the Backus \& Driffill (1985) thing can still apply. 
    \end{itemize}
\end{itemize}

Some smaller issues:
\begin{itemize}
    \item NT may have a conflict of interest --- it does not just care about prices being close to fundamentals. It might just want prices to have a high level, and also to reduce crashes (but care less about bubbles). 
\end{itemize}

\subsection{New sketch?}

Just occurred to me: perhaps the govt would like to have a dynamic plan? I can start with static, of course. 

TODO: the longer-term demand thing, what is a more appropriate model? Funding/intermediary? It might also be good to just list out the various things that can be tweaked. (Will Cong discussion)

\clearpage
\bibliographystyle{rfs}
\bibliography{refs}

\processdelayedfloats

\newpage
\appendix

    
 % Resetting for the actual figures and tables
\setcounter{figure}{0}
\setcounter{table}{0}
\renewcommand{\thetable}{\Alph{section}\arabic{table}}
\renewcommand{\thefigure}{\Alph{section}\arabic{figure}}

% \section{Additional Results}
% \label{app:additional}


\processdelayedfloats

\end{document}
